\section{哈代域中函数的比较}

命题 1. 同一哈代域中的两个函数可比较到任意阶 (V, p.~232)。

事实上，若 \(f\) 属于一个哈代域 \(\mathfrak{K}\) ，则对每个整数 \(n > 0\) 存在一个区间 \([x_0, +\infty[\) ， \(f\) 在其上 \(n\) 次可导，且其 \(n^{th}\) 导数在该区间上等于 \(\mathfrak{K}\) 中的一个函数。因此只需证明任意两个函数 \(f, g\) （属于 \(\mathfrak{K}\) ）可比较。若两函数之一在 \(+\infty\) 的某个邻域上恒为零，这是显然的；于是可限于两者都在 \(+\infty\) 的某个邻域上严格正的情形。但这时，对每个实数 \(t\) ， \(f - tg\) 等于 \(\mathfrak{K}\) 中的一个函数（在 \(+\infty\) 的某个邻域上），故在 \(+\infty\) 的某个邻域上具有常号，这就证明了命题 (V, p.~217, prop. 9)。

由这个命题立即推出：若哈代域 \(\mathfrak{K}\) 包含实常数（下文中我们总作此假定），且 \(f\) 与 \(g\) 是 \(\mathfrak{K}\) 中任意两个函数，则函数 \(e^f\) , \(e^g\) , \(\log |f|\) , \(\log |g|\) , \(|f|^{\alpha}\) , \(|g|^{\alpha}\) （ \(\alpha\) 为任意实数）、 \(\int_{a} f\) , \(\int_{a} g\) （ \(a\) 是某个区间 \([x_0, +\infty[\) 中的任意实数， \(f\) 与 \(g\) 在其上为规则函数）中的任意两个都可比较（当它们有定义时）；事实上，这些函数中的任意两个都属于把它们相继添加到 \(\mathfrak{K}\) 而得到的某个哈代域。

同样，每个函数 \(f(x)\) （属于哈代域 \(\mathfrak{K}\) ）都与 \(x\) 可比较，因为 \(x\) 与 \(f(x)\) 都属于把 \(x\) 添加到 \(\mathfrak{K}\) 而得到的哈代域。于是（特别地）可以断定， \(f\) 与每个幂 \(x^{\alpha}\) 以及与 \(\log x\) 、与 \(e^{x}\) 都可比较到任意阶。

我们还看到，若 \(f\) 与 \(g\) 属于同一哈代域 \(\mathfrak{K}\) ， \(g(x) > 0\) 于某个区间 \([x_0, +\infty[\) 上，且 \(g(x)\) 趋于 0 或 \(+\infty\) （当 \(x\) 趋于 \(+\infty\) 时），则 \(f\) 相对于 \(g\) 的阶 (V, p.~219) 总有定义。

因此，V, p.~233 的命题 8 可应用于哈代域中的每个函数 \(f\) ，并给出：

\(1^{\circ}\) 若 \(f\) 的阶为 \(+\infty\) （相对于 \(x\) ），则 \(\int_{a}^{x} f(t) dt \sim (f(x))^{2} / f'(x)\) 。

\(2^{\circ}\) 若 \(f\) 的阶为 \(\mu > -1\) （相对于 \(x\) ），则 \(\int_{a}^{x} f(t) dt \sim \frac{1}{\mu + 1} xf(x)\) 。

\(3^{\circ}\) 若 \(f\) 的阶为 \(\mu < -1\) （相对于 \(x\) ），则 \(\int_{x}^{+\infty} f(t) dt \sim -\frac{1}{\mu + 1} xf(x)\) 。

\(4^{\circ}\) 若 \(f\) 的阶为 \(-\infty\) （相对于 \(x\) ），则 \(\int_{x}^{+\infty} f(t) dt \sim -(f(x))^{2} / f'(x)\) 。

此外，我们有下列命题：

命题 2. 设 \(f\) 是属于哈代域 \(\mathfrak{K}\) 的一个函数。

\(1^{\circ}\) 若 \(f\) 相对于 \(x\) 的阶为无穷，则对每个整数 \(n > 0\) ，

\[
f ^ {(n)} (x) \sim \frac {(f ^ {\prime} (x)) ^ {n}}{(f (x)) ^ {n - 1}}.\tag{2}
\]

\(2^{\circ}\) 若 \(f\) 的阶为有限数 \(\mu\) （相对于 \(x\) ），则对每个 \(n > 0\) ，

\[
\begin{array}{c} f ^ {(n)} (x) \sim \mu (\mu - 1) \ldots (\mu - n + 1) \frac {f (x)}{x ^ {n}} \\ \sim \frac {(\mu - 1) \ldots (\mu - n + 1)}{\mu^ {n - 1}} \frac {(f ^ {\prime} (x)) ^ {n}}{(f (x)) ^ {n - 1}} \end{array}\tag{3}
\]

除非 \(\mu\) 是一个整数 \(\geqslant 0\) 且 \(n > \mu\) 。

\(1^{\circ}\) 若 \(f\) 相对于 \(x\) 的阶为无穷，则有 \(\log |f| \gg \log x\) ，于是，由于 \(\log |f|\) 与 \(\log x\) 可比较到任意阶，有 \(f'/f \gg 1/x\) 。令 \(g = f'/f\) ；因为 \(g\) 等于 \(\mathfrak{K}\) 中的一个函数（在 \(+\infty\) 的某个邻域上），由 \(1/g \ll x\) 推出 \(g'/g^{2} \ll 1\) ，从而 \(g'/g \ll g = f'/f\) ，或者说 \(fg' \ll gf'\) 。由关系式 \(f' = fg\) 求导得到

\[
f ^ {\prime \prime} = f g ^ {\prime} + g f ^ {\prime} \sim g f ^ {\prime}
\]

或者说 \(f'' / f' \sim f' / f\) 。把同样的论证应用于 \(f^{(n)}\) 以代替 \(f\) ，并对 \(n\) 作归纳，便得 \(f^{(n)} / f^{(n-1)} \sim f' / f\) ；由此得关系式 (2)。

\(2^{\circ}\) 若 \(f\) 的阶为有限数 \(\mu\) （相对于 \(x\) ），且 \(\mu \neq 0\) ，则有 \(\log |f| \sim \mu \log x\) ，求导便得 \(f'(x) \sim \mu \frac{f(x)}{x}\) ；由此推出 \(f'\) 的阶为 \(\mu - 1\) （相对于 \(x\) ），这使我们可以对 \(n\) 作归纳（只要 \(\mu \neq n\) ）而应用同样的论证，从而当 \(\mu\) 不是整数 \(\geqslant 0\) 且 \(< n\) 时得到公式 (3)。

当 \(f\) 的阶为整数 \(p \geqslant 0\) （相对于 \(x\) ）时，可写 \(f(x) = x^p f_1(x)\) ，其中 \(f_1\) 相对于 \(x\) 的阶为 0。由命题 2 有

\[
f ^ {(p)} \sim p! f _ {1}.
\]

为求出阶数 \(n > p\) 的导数，可限于 \(p = 0\) 的情形。这时有 \(\log |f| \ll \log x\) ，从而 \(f'(x) / f(x) \ll 1 / x\) ，换言之 \(xf'(x) \ll f(x)\) ；若 \(f\) 不等价于一个常数 \(k \neq 0\) ，对此关系式求导 (V, p.~232, prop. 7) 便得 \(xf''(x) + f'(x) \ll f'(x)\) ，这蕴含 \(xf''(x) \sim -f'(x)\) 。利用这个公式，对 \(n\) 作归纳可见 \(f^{(n)}\) 的阶 \(\leqslant -n\) （相对于 \(x\) ），并且

\[
f ^ {(n)} \sim (- 1) ^ {n + 1} (n - 1)! \frac {f ^ {\prime} (x)}{x ^ {n - 1}}.\tag{4}
\]

若 \(f\) 等价于一个常数 \(k \neq 0\) ，则有 \(f(x) = k + f_2(x)\) ，其中 \(f_2 \ll 1\) ，于是归结为研究 \(f_2\) 的导数。

